\scanpage{73}
We are going to prove the auxiliary proposition

\uline{Proposition 3} We have, for $u\in L_0(E,F)$, and $0\leq m\leq n$

\[
\begin{array}{@{}l@{\qquad}l@{}}
\textcircled{19}\ \deleted{\textcircled{18}} &
\displaystyle \sigma_n^m(u)=\operatorname*{Sup}_{\left\{\substack{v_n\in L(F,E)\\\lVert v_n\rVert\leq1\\\operatorname{rang}v_n\leq n}\right.}|\alpha_m(v_nu)|=\operatorname*{Sup}_{v_n}\{\alpha_{\unsure{n|m}}(|v_nu|)
\end{array}
\]

\deleted{\unsure{Indeed}}.

Indeed, let us prove, \deleted{that $|\alpha_m(v_nu)|=|\operatorname{Tr}\overset{m}{\Lambda}v_nu|\leq$} \deleted{It suffices} to prove the inequality \deleted{$|\operatorname{Tr}(\overset{m}{\Lambda}v_n)(\overset{m}{\Lambda}u)|\leq\lVert\overset{m}{\Lambda}\unsure{v_n}\rVert\,\lVert\overset{m}{\Lambda}u\rVert_1\leq\overset{m}{\Lambda}u$} $|\alpha_m(v_nu)|\leq\underline{\sigma_n^m}(u)$, since we have $\rho_i(v_nu)\leq\rho_i(u)$ for every $i$ (\unsure{---}), it suffices to establish $|\alpha_m(v_nu)|\leq\sigma_n^m(v_nu)$, hence the

\uline{lemma} If $u\in L(E,E)$ is of rank $\leq n$, we have $|\alpha_m(u)|\leq\{\alpha_m(|u|)\}\leq\sigma_n^m(u)$. Now we have, if $w$ is the restriction of $u$ to $u(E)$, \deleted{and $\alpha$ the identity map $u(E)\longrightarrow E$;}

\[
|\alpha_m(\deleted{\unsure{uv}}u)|=|\alpha_m(w)|\leq\deleted{\unsure{\alpha_m(|v|)}}\sum_{1\leq i_1\mathord{-}\leq i_m\leq n}\rho_{i_1}(w)\mathord{-}\rho_{i_m}(w),
\]

and since \deleted{\unsure{that}} $w=\alpha^*u\alpha$ ($\alpha$ identity map $u(E)\longrightarrow E$) whence $\rho_i(w)\leq\lVert\alpha^*\rVert\lVert\alpha\rVert\rho_i(u)\leq\rho_i(u)$, the desired formula.

\margin{\unsure{1|$\uparrow$}}

\deleted{\unsure{Let us show}} Let us show that the bound in (\unsure{19|18}) is \uuline{attained} : we may suppose $E$ of dimension $\geq n$, \deleted{\unsure{then}} and if \deleted{\unsure{one is} $\sum_{\unsure{i=1|i}}^{\infty}$ \unsure{of the} \unsure{$\rho_i e_i\otimes f_i$}}, $u=\sum\rho_i\bar e_i\otimes f_{\unsure{2|i}}$, we take $v_n=\sum_1^n\bar f_i\otimes e_i$, whence $v_nu=\sum_1^n\rho_i\bar e_i\otimes e_i$, and $\alpha_m(v_nu)=S_n^m(u)$ (remark : if $F=E$, and $u$ Hermitian, then we can take $v$ a Hermitian projector of rank $n$ for $E$ of dimension $\leq n$).
Let us prove, now (16) : \deleted{\unsure{Now} (17) reduces us to proving that} \inserted{(18) reduces us to proving} that $|\alpha_m(w(u+v))|\deleted{\unsure{\leq\sum_{p=0}^{n}\alpha_p(wu)}}$ is bounded above by the second side for \deleted{\unsure{every}} $w\in L(F,E)$ such that $\lVert w\rVert\leq1$, $\operatorname{rang}w\leq n$, \deleted{\unsure{now it is}} \deleted{\unsure{we}} we have ((14))

\[
|\alpha_{\unsure{n|m}}(w(u+v))|\leq\alpha_{\unsure{n|m}}(|wu+wv|)\leq\sum_{p=0}^{m}\alpha_p(|wu|)\alpha_{m-p}(|wv|)
\]

\deleted{on the other hand we have}

\[
\deleted{\alpha_p(|wu|)=\lVert\overset{p}{\Lambda}wu\rVert_1=\lVert(\overset{p}{\Lambda}w)(\overset{p}{\Lambda}u)\rVert_1\leq\lVert\overset{\unsure{n|p}}{\Lambda}w\rVert\,\lVert\overset{p}{\Lambda}u\rVert_1\leq\lVert\overset{\unsure{n|p}}{\Lambda}u\rVert_1=\unsure{\alpha_p(|u|)}}
\]

\deleted{now the operation \unsure{of the}} Now (\unsure{19|18}) shows that $\alpha_p(|wu|)\leq\sigma_n^p(u)$ and $\alpha_{m-p}(|wv|)\leq\sigma_n^{m-p}(v)$, whence the formula (18),
\margin{74}
